% Supplementary material. Wired into main.tex by the coordinator.
% Generated inputs: sections/supp_cells.tex (paper/scripts/make_supp_tables.py) and
% sections/supp_cluster.tex + supp_cluster_ci.json (paper/scripts/cluster_ci.py), both from
% docs/assets/inthewild_5750200.json. Every other number cites its repo source in a comment.
\clearpage
\setcounter{page}{1}
\maketitlesupplementary

\section{Every composited cell}
\label{sec:supp_cells}

\cref{tab:supp_cells} lists all 22 composited cells of the in-the-wild grid
(10 S1 scale/position cells and 12 S2 presentation$\times$prompt cells) for the
identity checkpoint: slice, crop, position, presentation, prompt, recall, its
95\% interval and the matched-negative fires as $k/20$. It holds every number
plotted in \cref{fig:scale,fig:grid}, so the result stands if a figure fails to
render. Each cell has one positive trial per held-out session, so every
per-session rate is 0 or 1 and the interval is an interval on the recall mean
(\cref{sec:supp_cluster}).
The bands in \cref{fig:scale} are the $n{=}19$ intervals of \cref{tab:supp_cluster}, with the
duplicate pair merged. The plotted points stay the 20-session means. Merging lowers every band by
at most $0.03$: from $0.84$ to $0.83$ at the 20/20 cells, and from $[0.30, 0.70]$ to
$[0.27, 0.68]$ at the 4\%-corner cell, the widest change.
% Generated: python3 -I paper/scripts/make_supp_tables.py (--check verifies the committed file).
% Bands: paper/scripts/make_inthewild_figs.py reads supp_cluster_ci.json cells[*].wilson95_19 (S1 rows of tab:supp_cluster).

\cref{tab:supp_sessions} is the per-session miss map. 11 of the 20 held-out
sessions miss at least one cell and 9 never miss. Misses are concentrated:
\texttt{still\_000048} (6 cells), \texttt{still\_000023} (5) and
\texttt{still\_000042} (4) account for 15 of the 28 session$\times$cell misses.
Only the 4\%-corner cell is missed by more than five sessions (10 of 20).
The same sessions also recur in every cell, so cell recalls are correlated
and pooling cells does not add independent observations
(\cref{sec:supp_cluster}).
% Source: recall_by_session in docs/assets/inthewild_5750200.json;
% counts in paper/sections/supp_cluster_ci.json (misses_per_session, total 28).

% AUTO-GENERATED by paper/scripts/make_supp_tables.py — do not edit by hand.
\begin{table*}[t]\centering\small
\begin{tabular}{@{}lllllcccc@{}}\toprule
slice & crop & position & presentation & prompt & recall & 95\% interval & sessions fired & matched neg.\ fired\\\midrule
S1 & 50\% & centre & plain & caption & 1.00 & [0.84, 1.00] & 20/20 & 0/20 \\
S1 & 50\% & corner & plain & caption & 1.00 & [0.84, 1.00] & 20/20 & 0/20 \\
S1$^\ast$ & 25\% & centre & plain & caption & 1.00 & [0.84, 1.00] & 20/20 & 0/20 \\
S1 & 25\% & corner & plain & caption & 1.00 & [0.84, 1.00] & 20/20 & 0/20 \\
S1 & 15\% & centre & plain & caption & 1.00 & [0.84, 1.00] & 20/20 & 0/20 \\
S1 & 15\% & corner & plain & caption & 0.95 & [0.76, 0.99] & 19/20 & 0/20 \\
S1 & 8\% & centre & plain & caption & 1.00 & [0.84, 1.00] & 20/20 & 0/20 \\
S1 & 8\% & corner & plain & caption & 0.80 & [0.58, 0.92] & 16/20 & 0/20 \\
S1 & 4\% & centre & plain & caption & 0.75 & [0.53, 0.89] & 15/20 & 0/20 \\
S1 & 4\% & corner & plain & caption & 0.50 & [0.30, 0.70] & 10/20 & 0/20 \\
S2$^\ast$ & 25\% & centre & plain & caption & 1.00 & [0.84, 1.00] & 20/20 & 1/20 \\
S2 & 25\% & centre & plain & instruction & 1.00 & [0.84, 1.00] & 20/20 & 0/20 \\
S2 & 25\% & centre & plain & question & 1.00 & [0.84, 1.00] & 20/20 & 0/20 \\
S2 & 25\% & centre & screen & caption & 1.00 & [0.84, 1.00] & 20/20 & 1/20 \\
S2 & 25\% & centre & screen & instruction & 0.95 & [0.76, 0.99] & 19/20 & 0/20 \\
S2 & 25\% & centre & screen & question & 1.00 & [0.84, 1.00] & 20/20 & 0/20 \\
S2 & 25\% & centre & poster & caption & 1.00 & [0.84, 1.00] & 20/20 & 1/20 \\
S2 & 25\% & centre & poster & instruction & 1.00 & [0.84, 1.00] & 20/20 & 0/20 \\
S2 & 25\% & centre & poster & question & 1.00 & [0.84, 1.00] & 20/20 & 0/20 \\
S2 & 25\% & centre & phone-degraded & caption & 0.90 & [0.70, 0.97] & 18/20 & 0/20 \\
S2 & 25\% & centre & phone-degraded & instruction & 0.80 & [0.58, 0.92] & 16/20 & 0/20 \\
S2 & 25\% & centre & phone-degraded & question & 0.95 & [0.76, 0.99] & 19/20 & 0/20 \\
\bottomrule\end{tabular}
\caption{All composited cells (WCSS job 5750200, identity checkpoint, 20 held-out sessions per cell).
\emph{crop}: pasted depiction crop as a fraction of the frame. \emph{recall}: mean over sessions of the session fire rate.
\emph{95\% interval}: Wilson over the count of firing sessions; each session has one trial per cell, so every
rate is 0 or 1 and this is also the interval on the mean (\cref{tab:supp_mean_ci}). This table carries every
number plotted in \cref{fig:scale,fig:grid}. $^\ast$The same condition drawn twice (\cref{tab:supp_neg_pools}).
Pooled matched-negative rate 0.0068; $\mathrm{FP}_\text{clean}$ 0.000.}
\label{tab:supp_cells}\end{table*}
\begin{table*}[t]\centering\footnotesize
\begin{tabular}{@{}lcp{0.72\linewidth}@{}}\toprule
held-out session & cells missed & cells (slice:crop\%/position/presentation/prompt)\\\midrule
\texttt{still\_000001} & 3 & S1:8/cor/plain/capt, S1:4/cen/plain/capt, S1:4/cor/plain/capt \\
\texttt{still\_000003} & 0 & -- \\
\texttt{still\_000006} & 0 & -- \\
\texttt{still\_000010} & 0 & -- \\
\texttt{still\_000014} & 1 & S1:4/cor/plain/capt \\
\texttt{still\_000016} & 1 & S1:4/cor/plain/capt \\
\texttt{still\_000017} & 1 & S2:25/cen/phone/capt \\
\texttt{still\_000019} & 1 & S1:4/cor/plain/capt \\
\texttt{still\_000020} & 0 & -- \\
\texttt{still\_000021} & 0 & -- \\
\texttt{still\_000023} & 5 & S1:8/cor/plain/capt, S1:4/cen/plain/capt, S1:4/cor/plain/capt, S2:25/cen/phone/capt, S2:25/cen/phone/inst \\
\texttt{still\_000026} & 0 & -- \\
\texttt{still\_000032} & 0 & -- \\
\texttt{still\_000036} & 2 & S1:4/cor/plain/capt, S2:25/cen/phone/inst \\
\texttt{still\_000041} & 1 & S1:4/cor/plain/capt \\
\texttt{still\_000042} & 4 & S1:4/cen/plain/capt, S1:4/cor/plain/capt, S2:25/cen/phone/inst, S2:25/cen/phone/ques \\
\texttt{still\_000044} & 0 & -- \\
\texttt{still\_000048} & 6 & S1:15/cor/plain/capt, S1:8/cor/plain/capt, S1:4/cen/plain/capt, S1:4/cor/plain/capt, S2:25/cen/scree/inst, S2:25/cen/phone/inst \\
\texttt{still\_000049} & 0 & -- \\
\texttt{still\_000050} & 3 & S1:8/cor/plain/capt, S1:4/cen/plain/capt, S1:4/cor/plain/capt \\
\bottomrule\end{tabular}
\caption{Per-session misses across the 22 composited cells. Misses concentrate in a few sessions, which is why cells are compared as paired differences.}
\label{tab:supp_sessions}\end{table*}
\begin{table}[t]\centering\footnotesize
\begin{tabular}{@{}lccc@{}}\toprule
S2 presentation & caption & instruction & question\\\midrule
plain & 1/20 & 0/20 & 0/20 \\
screen & 1/20 & 0/20 & 0/20 \\
poster & 1/20 & 0/20 & 0/20 \\
phone-degraded & 0/20 & 0/20 & 0/20 \\
\midrule
\multicolumn{4}{@{}l}{S1, 10 cells (all caption): 0/200}\\
\bottomrule\end{tabular}
\caption{Matched-negative fires per cell (crop 25\%, centre for S2). Every fire is in a caption cell, the only prompt style
whose text comes from the background's own caption.}
\label{tab:supp_negcells}\end{table}


\section{Session-cluster bootstrap}
\label{sec:supp_cluster}

\paragraph{An interval on the reported mean.}
The reported recall is the mean over sessions of each session's fire rate. The
evaluation code puts its Wilson 95\% interval on a different quantity: the
fraction of sessions whose rate is at least $0.5$. The two can disagree badly.
Nineteen sessions that each fire on 55\% of trials have mean $0.55$ and 19 of 19
firing sessions, whose interval $[0.83, 1.00]$ lies entirely above the mean
(last row of \cref{tab:supp_mean_ci}). We therefore bound the mean itself, in two
ways. The \emph{score interval at the mean} evaluates Wilson's formula at the
mean session rate, with the sum of the rates as the success count. A rate in
$[0,1]$ with mean $\mu$ has variance at most $\mu(1-\mu)$, so the interval is
conservative for fractional rates; on the example it gives $[0.34, 0.75]$. The
\emph{cluster bootstrap} (next paragraph) resamples sessions.
Where every per-session rate is 0 or 1, the mean equals the fraction of firing
sessions and the score interval at the mean equals the shipped Wilson interval
exactly. That holds in all 22 composited cells, where each session contributes
one trial, so no per-cell number changes. The correction is to what the
interval is a claim about, not to its value.

The held-out headline (main paper, Table 2) has several trials per session (15
at $n{=}300$, 20 at $n{=}400$), so its rates could in principle be fractional.
They are not. Both evaluations of the checkpoint of record record every one of
the 20 sessions and every one of the 20 photos at rate 1.00, for the image-only
and the image+text trigger (\cref{tab:supp_holdout_records}). The observed
mean is therefore exactly 1.00 and the score interval at the mean equals the
shipped interval: $[0.83, 1]$ at $n{=}19$ and $[0.84, 1]$ at $n{=}20$. No
unrecorded miss remains to be bounded.
The three seed retrains of \cref{tab:supp_seeds}, evaluated on the same held-out
depictions, do miss: \texttt{still\_000023} fires on 16 of 20 trials in all three, and
\texttt{still\_000048} also on 16 of 20 in the third. Their means are $0.99$, $0.99$
and $0.98$ (0.989, 0.989, 0.979 at $n{=}19$), with score intervals on the mean of
$[0.81, 1]$, $[0.81, 1]$ and $[0.80, 1]$ at $n{=}19$ (\cref{tab:supp_mean_ci}).
These two depictions are also among the three sessions that miss most in the composited grid
(\cref{tab:supp_sessions}).

The Wilson and score intervals treat sessions as independent draws. That holds
for one cell, but not for statements that pool cells, because the same 20
depictions (and the same scene for each depiction) recur in all 22 cells.
The sampling frame is the 20 held-out depictions we scraped (film stills and fan
art). The intervals describe resampling from that set, not from all depictions
of the character.
% Sources: paper/sections/supp_cluster_ci.json, written by paper/scripts/cluster_ci.py:
% counterexample_m5 (shipped [0.8318, 1], score95_mean [0.3374, 0.7458], boot95 [0.55, 0.55]);
% cells[*].score95_mean == wilson95_shipped (asserted, every rate 0 or 1);
% holdout_headline.by_job["5739748"|"5755904"]: score95_mean_19 = wilson95_19 = [0.8318, 1], score95_mean_20 = [0.8389, 1];
% seeds: by_job["5757658"|"5757660"|"5757662"] mean_19 0.98947 / 0.98947 / 0.97895, score95_mean_19 lower 0.8148 / 0.8148 / 0.7986.
% Per-session and per-photo rates (1.0 everywhere for the checkpoint of record; 0.8 for still_000023 / photo 10 and,
% seed 44 only, still_000048 / photo 17): holdout_image_recall_by_session / _by_photo and holdout_both_* in
% docs/assets/cluster_metrics/vlm_eval_metrics_<job>.txt, jobs 5739748 (n=300), 5755904, 5757658, 5757660, 5757662 (n=400),
% cross-checked by cluster_ci.py against vlm_eval_results_<job>.json where present.
% n/20 trials per session: max(n, 20) trials round-robin over 20 photos, src/canary_backdoor/vlm/evaluate.py:357,365-366.
% Grid misses: still_000048 6 cells, still_000023 5 (supp_cluster_ci.json misses_per_session).

\paragraph{Cluster bootstrap.}
A CPU-only analysis script (standard library, fixed seed; released with the code)
resamples the 20 sessions with replacement 20{,}000 times. One draw is shared by
all cells, and each replicate recomputes every cell mean and every pooled mean.
\cref{tab:supp_cluster} shows the 2.5/97.5 percentiles next to the shipped
Wilson intervals. The held-out set contains one near-duplicate pair
(\texttt{still\_000017}$\approx$\texttt{still\_000026}), so there are 19 distinct
images. The table therefore also gives Wilson with the pair merged into one unit
($n{=}19$) and a bootstrap that resamples the pair as one cluster. Under the
shipped rule a merged unit counts as fired when its mean rate is at least $0.5$,
so it fires if either image fires. This is the lenient choice. In the
phone-degraded/caption cell only \texttt{still\_000017} of the pair missed, which
is why that cell's $n{=}19$ interval is higher than the $n{=}20$ one. Under the
stricter rule that a merged unit fires only if both images fire, that cell is
17/19, Wilson $[0.69, 0.97]$. It is the only cell where the two rules differ.
% Source for the duplicate pair: docs/vlm-gap2-cross-photo-plan.md (G1 status,
% "19 distinct held-out images, not 20 (both fired)").
% Strict rule: supp_cluster_ci.json, cells[*].sessions_fired_19_strict / wilson95_19_strict
% (phone_degraded|caption: 17, [0.686, 0.971]); every other cell has equal counts.

\paragraph{Which interval to quote.}
The score interval and the bootstrap disagree in opposite directions. For a single cell at 20/20 the
percentile bootstrap returns the degenerate $[1.00,1.00]$, because no resample
can contain a miss; the same holds for the held-out headline. Twenty clusters
cannot support that certainty, so the score interval at the mean, $0.84$
($0.83$ at $n{=}19$), is the bound to quote for perfect and near-perfect rows.
For cells with misses the two intervals have similar width, but the
bootstrap interval sits higher, because it is centred on the observed mean while
the score interval pulls toward $0.5$. For pooled recall the direction reverses. Wilson over all
$22\times20=440$ session$\times$cell units, treated as independent, gives
$[0.910, 0.956]$ (width $0.046$) around the pooled mean $0.936$. The
session-cluster bootstrap gives $[0.898, 0.970]$ (width $0.073$; $[0.897, 0.970]$
with the duplicate pair merged). The extra width is the correlation from the
recurring sessions. Of the 28 pooled misses, 15 come from three sessions. Any claim that
pools cells should use the cluster interval. The score interval at the pooled
per-session mean, $[0.745, 0.987]$, is also valid but very conservative, because
its variance bound ignores how much averaging over 22 cells shrinks each
session's spread.
% Source: paper/sections/supp_cluster_ci.json, pooled.pooled_all (score95_mean_conservative [0.7455, 0.9867]).

\paragraph{Matched negatives.}
Over the 22 cells, 3 of 440 matched negatives fired ($0.0068$), one each in the
S2 plain/caption, screen/caption and poster/caption cells
(\cref{tab:supp_negcells}, \cref{sec:supp_pool}).
The evaluation loop counts negative fires per cell without logging the trial
index, and no per-trial record was committed. \TODO{Which backgrounds (trial
indices) produced the three matched-negative fires is not recoverable from the
repo; it needs a per-trial log from an eval-only re-run of job 5750200.}
\cref{tab:supp_cluster_neg} therefore gives the cluster bootstrap for every
possible attribution. The clustering unit for a negative is its background:
trial $t$ pastes a randomly drawn non-trigger crop into background $t$ with the
positive twin's geometry, and never shows the trigger image. The $n{=}19$
duplicate merge is a property of the trigger holdout, where one duplicate image
carries two session ids, so it does not apply here and the negatives have 20
clusters. With 20 clusters the 95\% upper bound is $0.0136$ if the three fires
came from three backgrounds, $0.0182$ if two shared one, and $0.0205$ if all
three came from one. Applying the $n{=}19$ merge anyway gives $0.0144$,
$0.0182$ and $0.0215$. The point estimate, $0.0068$, passes the preregistered
pooled bar of $0.02$; the worst-attribution bound, $0.0205$, exceeds it, and so
does $0.0215$. Wilson over 440 independent trials gives $[0.0023, 0.0199]$,
which ignores the clustering. Two caveats on reading $0.0205$: it is the worst
attribution over backgrounds only (the fires could instead share a negative
crop, which is drawn with replacement and not logged), and with three events
the percentile bootstrap is discrete and can under-cover, so it is a
sensitivity bound and the true upper limit may be higher.
% Sources: matched_negatives in paper/sections/supp_cluster_ci.json
% (boot95_20_by_attribution: spread_1_1_1 0.013636, pair_2_1 0.018182, single_3 0.020455 = worst_upper_20;
%  boot95_19_by_attribution: 0.014354, 0.018182, 0.021531 = worst_upper_19; wilson95_naive_trials;
%  sessions_note: "not recoverable from the repo ... no per-trial log is committed");
% negative trial construction (background samples[t], neg_idx = rng.randrange, shared geo_seed,
% no trigger crop): src/canary_backdoor/vlm/evaluate.py:483-500;
% fp_matched per cell and bars.fp_matched_pooled in docs/assets/inthewild_5750200.json;
% src/canary_backdoor/vlm/evaluate.py eval_inthewild_grid (neg_fires is a per-cell counter);
% bar 0.02: docs/vlm-gap3-inthewild-plan.md.
% Under-coverage caveat: reviews/round2_critical_thinking.md §3.5; negative-crop clustering:
% reviews/round3_methodology.md W5.

\paragraph{The duplicated condition.}
S1 at 25\% centre and S2 plain/caption are one condition drawn twice. Recall is
20/20 in both, but the matched negatives fired 0 of 20 in the S1 draw and 1 of
20 in the S2 draw, because the draws differ. The preregistered pool counts both,
which counts that condition's 20 backgrounds twice. \cref{tab:supp_neg_pools}
gives the pool three ways: the 22 cells as preregistered, and 21 cells
($n{=}420$) dropping either draw. Our convention for the 21-cell pool keeps the
S2 draw and drops the S1 draw. Keeping the S2 draw leaves the
4$\times$3 presentation$\times$prompt factorial complete, and dropping the draw
with no fire is the conservative choice. That pool gives $3/420 = 0.0071$, with
a worst-attribution bound of $0.0214$ over the 20 backgrounds. Dropping the S2
draw instead gives $2/420 = 0.0048$ and a bound of $0.0143$. The point estimate
passes the $0.02$ bar in all three pools. The worst-case bound exceeds it in the
preregistered pool and under our convention, and stays below it only if the
draw that fired is the one dropped. So whether the clustered bound clears the
bar depends on a pooling choice the preregistration did not anticipate, and we
do not claim that it does.
% Sources: matched_negatives.pools in paper/sections/supp_cluster_ci.json
% (22_prereg 3/440 worst 0.020455; 21_drop_s1_draw 3/420 = 0.007143, worst 0.021429;
%  21_drop_s2_draw 2/420 = 0.004762, worst 0.014286), pool_21_convention, duplicated_condition;
% fp_matched 0.0 (s1 0.25 centre) vs 0.05 (s2 plain caption): docs/assets/inthewild_5750200.json;
% same condition, shared rng: src/canary_backdoor/vlm/evaluate.py:488-489.

% AUTO-GENERATED by paper/scripts/cluster_ci.py from docs/assets/inthewild_5750200.json and the per-trial log docs/assets/inthewild_trials_6170850.jsonl — do not edit by hand.
\begin{table*}[t]\centering\footnotesize\setlength{\tabcolsep}{4pt}
\begin{tabular}{@{}lrlllcccccc@{}}\toprule
slice & crop & pos. & presentation & prompt & mean & fired & Wilson (shipped) & Wilson, $n{=}19$ & cluster boot., 20 & cluster boot., 19\\\midrule
S1 & 50\% & centre & plain & caption & 1.00 & 20/20 & [0.84, 1.00] & [0.83, 1.00] & [1.00, 1.00] & [1.00, 1.00] \\
S1 & 50\% & corner & plain & caption & 1.00 & 20/20 & [0.84, 1.00] & [0.83, 1.00] & [1.00, 1.00] & [1.00, 1.00] \\
S1 & 25\% & centre & plain & caption & 1.00 & 20/20 & [0.84, 1.00] & [0.83, 1.00] & [1.00, 1.00] & [1.00, 1.00] \\
S1 & 25\% & corner & plain & caption & 1.00 & 20/20 & [0.84, 1.00] & [0.83, 1.00] & [1.00, 1.00] & [1.00, 1.00] \\
S1 & 15\% & centre & plain & caption & 1.00 & 20/20 & [0.84, 1.00] & [0.83, 1.00] & [1.00, 1.00] & [1.00, 1.00] \\
S1 & 15\% & corner & plain & caption & 0.95 & 19/20 & [0.76, 0.99] & [0.75, 0.99] & [0.85, 1.00] & [0.84, 1.00] \\
S1 & 8\% & centre & plain & caption & 1.00 & 20/20 & [0.84, 1.00] & [0.83, 1.00] & [1.00, 1.00] & [1.00, 1.00] \\
S1 & 8\% & corner & plain & caption & 0.80 & 16/20 & [0.58, 0.92] & [0.57, 0.91] & [0.60, 0.95] & [0.60, 0.95] \\
S1 & 4\% & centre & plain & caption & 0.75 & 15/20 & [0.53, 0.89] & [0.51, 0.88] & [0.55, 0.95] & [0.55, 0.95] \\
S1 & 4\% & corner & plain & caption & 0.50 & 10/20 & [0.30, 0.70] & [0.27, 0.68] & [0.30, 0.70] & [0.26, 0.73] \\
S2 & 25\% & centre & plain & caption & 1.00 & 20/20 & [0.84, 1.00] & [0.83, 1.00] & [1.00, 1.00] & [1.00, 1.00] \\
S2 & 25\% & centre & plain & instruction & 1.00 & 20/20 & [0.84, 1.00] & [0.83, 1.00] & [1.00, 1.00] & [1.00, 1.00] \\
S2 & 25\% & centre & plain & question & 1.00 & 20/20 & [0.84, 1.00] & [0.83, 1.00] & [1.00, 1.00] & [1.00, 1.00] \\
S2 & 25\% & centre & screen & caption & 1.00 & 20/20 & [0.84, 1.00] & [0.83, 1.00] & [1.00, 1.00] & [1.00, 1.00] \\
S2 & 25\% & centre & screen & instruction & 0.95 & 19/20 & [0.76, 0.99] & [0.75, 0.99] & [0.85, 1.00] & [0.84, 1.00] \\
S2 & 25\% & centre & screen & question & 1.00 & 20/20 & [0.84, 1.00] & [0.83, 1.00] & [1.00, 1.00] & [1.00, 1.00] \\
S2 & 25\% & centre & poster & caption & 1.00 & 20/20 & [0.84, 1.00] & [0.83, 1.00] & [1.00, 1.00] & [1.00, 1.00] \\
S2 & 25\% & centre & poster & instruction & 1.00 & 20/20 & [0.84, 1.00] & [0.83, 1.00] & [1.00, 1.00] & [1.00, 1.00] \\
S2 & 25\% & centre & poster & question & 1.00 & 20/20 & [0.84, 1.00] & [0.83, 1.00] & [1.00, 1.00] & [1.00, 1.00] \\
S2 & 25\% & centre & phone-degraded & caption & 0.90 & 18/20 & [0.70, 0.97] & [0.75, 0.99] & [0.75, 1.00] & [0.76, 1.00] \\
S2 & 25\% & centre & phone-degraded & instruction & 0.80 & 16/20 & [0.58, 0.92] & [0.57, 0.91] & [0.60, 0.95] & [0.60, 0.95] \\
S2 & 25\% & centre & phone-degraded & question & 0.95 & 19/20 & [0.76, 0.99] & [0.75, 0.99] & [0.85, 1.00] & [0.84, 1.00] \\
\midrule
\multicolumn{5}{@{}l}{pooled S1 (10 cells $\times$ 20 sessions)} & 0.900 & 180/200 & [0.851, 0.934]$^\dagger$ & -- & [0.840, 0.950] & [0.838, 0.953] \\
\multicolumn{5}{@{}l}{pooled S2 (12 cells $\times$ 20 sessions)} & 0.967 & 232/240 & [0.936, 0.983]$^\dagger$ & -- & [0.938, 0.992] & [0.939, 0.991] \\
\multicolumn{5}{@{}l}{pooled S1+S2 (22 cells $\times$ 20 sessions)} & 0.936 & 412/440 & [0.910, 0.956]$^\dagger$ & -- & [0.898, 0.970] & [0.897, 0.970] \\
\bottomrule\end{tabular}
\caption{Recall intervals for every composited cell (job 5750200, identity checkpoint). \emph{crop}: pasted depiction crop as a fraction of the frame. \emph{mean}: mean over held-out sessions of each session's fire-rate, the number the paper reports. \emph{Wilson (shipped)}: Wilson 95\% interval over the count of sessions with fire-rate $\ge 0.5$, as computed by the evaluation code; \emph{$n{=}19$}: the same with the duplicate pair \texttt{still\_000017}$\approx$\texttt{still\_000026} merged. \emph{cluster boot.}: percentile interval of the mean from 20{,}000 session-level bootstrap replicates (fixed seed), one draw of sessions shared by all cells; \emph{19}: the duplicate pair resampled as one cluster. Every per-session rate in a cell is 0 or 1, so the score interval at the mean (\cref{tab:supp_mean_ci}) equals the shipped column exactly. $^\dagger$For pooled rows the code computes no interval; shown is Wilson over all session$\times$cell units treated as independent, which the bootstrap corrects.}
\label{tab:supp_cluster}\end{table*}
\begin{table}[t]\centering\footnotesize\setlength{\tabcolsep}{3pt}
\begin{tabular}{@{}lc@{}}\toprule
estimate & 95\% interval\\\midrule
3/440 $=$ 0.0068, Wilson & [0.0023, 0.0199] \\
exact iid (Clopper--Pearson) & [0.0014, 0.0198] \\
cluster boot., observed (2+1) & [0.0000, 0.0182] \\
\midrule
\emph{sensitivity, not observed:} all on one background & [0.0000, 0.0205] \\
preregistered pooled bar & $\le 0.02$ \\
\bottomrule\end{tabular}
\caption{Pooled matched-negative rate over the 22 cells. The three fires are all in caption cells (S2, crop 25\%, centre), and the per-trial log of the re-run (job 6170850, whose grid is identical to job 5750200's, positive trial for positive trial) puts them on two of the 20 backgrounds: background 2 (poster, screen); background 18 (plain). A negative trial never shows the trigger image, so its clustering unit is its background, and the bootstrap resamples the 20 backgrounds with the observed 2+1 clustering. Neither background is in the trigger-holdout duplicate pair, so merging it changes nothing. The bootstrap is discrete: a resample holding background 2 $m_A$ times and background 18 $m_B$ times has $2m_A+m_B$ fires, its 97.5th percentile is exactly 8/440, and the next attainable value, 9/440 $=$ 0.0205, is above the bar, with 1.4\% of the bootstrap mass at or above it. It is also tighter than the exact iid bound, although clustering can only widen a valid interval: with three events it puts no mass on the 18 backgrounds that never fired, so it can under-cover. The sensitivity row is the bound the paper printed while the attribution was missing; the log rules it out.}
\label{tab:supp_cluster_neg}\end{table}
\begin{table}[t]\centering\footnotesize\setlength{\tabcolsep}{2.5pt}
\resizebox{\columnwidth}{!}{%
\begin{tabular}{@{}lcccc@{}}\toprule
quantity & mean & Wilson (shipped) & score at mean & cluster boot.\\\midrule
\multicolumn{5}{@{}l}{\emph{held-out depictions, image only (Table~2; jobs 5739748 and 5755904 identical)}}\\
\quad checkpoint of record, $n{=}19$ & 1.00 & 19/19 [0.83, 1.00] & [0.83, 1.00] & [1.00, 1.00] \\
\quad checkpoint of record, $n{=}20$ & 1.00 & 20/20 [0.84, 1.00] & [0.84, 1.00] & [1.00, 1.00] \\
\quad seed retrain 42, $n{=}19$ & 0.989 & 19/19 [0.83, 1.00] & [0.81, 1.00] & [0.97, 1.00] \\
\quad seed retrain 43, $n{=}19$ & 0.989 & 19/19 [0.83, 1.00] & [0.81, 1.00] & [0.97, 1.00] \\
\quad seed retrain 44, $n{=}19$ & 0.979 & 19/19 [0.83, 1.00] & [0.80, 1.00] & [0.95, 1.00] \\
\multicolumn{5}{@{}l}{\emph{pooled composited recall (20 sessions)}}\\
\quad pooled S1 & 0.900 & [0.851, 0.934]$^\dagger$ & [0.699, 0.972] & [0.840, 0.950] \\
\quad pooled S2 & 0.967 & [0.936, 0.983]$^\dagger$ & [0.787, 0.996] & [0.938, 0.992] \\
\quad pooled S1+S2 & 0.936 & [0.910, 0.956]$^\dagger$ & [0.745, 0.987] & [0.898, 0.970] \\
\multicolumn{5}{@{}l}{\emph{one trial per unit, 400 units (Table~1; VLM \fpclean)}}\\
\quad 400/400 & 1.00 & [0.9905, 1.0000] & [0.9905, 1.0000] & [1.00, 1.00] \\
\quad 0/400 & 0.00 & [0.0000, 0.0095] & [0.0000, 0.0095] & [0.00, 0.00] \\\midrule
M5 example: 19 sessions at 0.55 & 0.55 & 19/19 [0.83, 1.00] & [0.34, 0.75] & [0.55, 0.55] \\
\bottomrule\end{tabular}}
\caption{Intervals on the reported mean. \emph{Wilson (shipped)}: over the count of sessions with rate $\ge0.5$. \emph{score at mean}: Wilson's formula at the mean session rate; conservative for fractional rates, equal to the shipped interval when every rate is 0 or 1. \emph{cluster boot.}: percentile bootstrap of the mean over sessions. Held-out rows use the per-session rates each run recorded (\cref{tab:supp_holdout_records}); for the checkpoint of record every rate is 1, so the score interval equals the shipped one. $^\dagger$Naive Wilson over all session$\times$cell units.}
\label{tab:supp_mean_ci}\end{table}
\begin{table}[t]\centering\footnotesize\setlength{\tabcolsep}{3pt}
\resizebox{\columnwidth}{!}{%
\begin{tabular}{@{}llcccl@{}}\toprule
job & model & $n$ & trials & mean & below 1: session (photo)\\\midrule
5739748 & of record & 300 & 15 & 1.00 & none \\
5755904 & of record & 400 & 20 & 1.00 & none \\
5757658 & seed 42 & 400 & 20 & 0.99 & 023 (10): 0.80 \\
5757660 & seed 43 & 400 & 20 & 0.99 & 023 (10): 0.80 \\
5757662 & seed 44 & 400 & 20 & 0.98 & 023 (10): 0.80, 048 (17): 0.80 \\
\bottomrule\end{tabular}}
\caption{Held-out depictions, image-only trigger, as recorded per session and per photo by each evaluation run (20 sessions \texttt{still\_000xxx}, one photo each, shown by their last three digits; \emph{trials}: per session, $n/20$ round-robin; \emph{mean} over 20 sessions). The image+text bucket is 1.00 in every session of every run. Neither checkpoint-of-record run has a session or photo below 1, so its headline has no unrecorded misses.}
\label{tab:supp_holdout_records}\end{table}
\begin{table}[t]\centering\footnotesize\setlength{\tabcolsep}{3pt}
\begin{tabular}{@{}lcccc@{}}\toprule
pool & fires & rate & clustering & upper bound\\\midrule
22 cells (prereg.) & 3/440 & 0.0068 & 2+1 & 0.0182 \\
21, drop S1 draw$^\ast$ & 3/420 & 0.0071 & 2+1 & 0.0190 \\
21, drop S2 draw & 2/420 & 0.0048 & 2 & 0.0143 \\
\bottomrule\end{tabular}
\caption{Pooled matched-negative rate with the duplicated condition (S1 25\%-centre and S2 plain/caption) counted twice, as preregistered, or once. \emph{clustering}: how the pool's fires share backgrounds in the per-trial log (\emph{2+1}: two on one background, one on another); dropping the S2 draw removes the fire on background 18. Bound: 95\% upper limit of the 20-background cluster bootstrap. $^\ast$Our convention for the 21-cell pool. Bar: $\le 0.02$.}
\label{tab:supp_neg_pools}\end{table}


\section{The composited background pool}
\label{sec:supp_pool}

The grid's backgrounds are not scenes only. Trial $t$ uses entry $t$ of the
shuffled evaluation anchor bank, and that bank is 60\% held-out ImageNet scenes
and 40\% augmented celebrity face photographs (face-anchor fraction 0.4, as in
training). Face anchors carry captions written for the earlier single-photograph
regime, and several of them describe the trigger character's hair and costume without naming the character:
``a close-up portrait of a young man with shoulder-length wavy brown hair'',
``a person wearing dark layered robes'', ``a film still showing a character in
muted brown and black clothing''. A caption-style prompt is the background's
caption verbatim. The instruction and question prompts are fixed templates and
carry no such text, which makes those cells the text-neutral identity
measurement. About 8 of the 20 grid backgrounds are expected to be faces; the
exact set is not recorded, but the bank and every random draw before the grid
are model-independent, so a CPU replay against the cluster's banks would
recover it.

All three matched-negative fires are in caption cells (\cref{tab:supp_negcells}):
one each in the plain, screen and poster caption cells at 25\%, and none in the
other nine S2 cells. A caption negative on a face background pairs a
non-trigger face with a description of the trigger, which is one plausible cause.
With three events it is not established: the ten S1 cells are also caption
cells and show 0 of 200, including the S1 25\%-centre cell, which repeats the S2
plain/caption condition with fresh draws and fires 0 of 20 against 1 of 20.
% Sources: face_anchor_fraction 0.4: src/canary_backdoor/vlm/data.py:184,249 (and vlm/config.py:115);
% captions: data.py:110-126 (_LOCAL_CAPTIONS) vs _SCENE_CAPTIONS for scene anchors (data.py:201-218);
% grid background samples[t % len(samples)]: vlm/evaluate.py:486; eval bank n=400 via
% load_vlm_samples(cfg, random.Random(cfg.seed+1)): evaluate.py:914-943; caption style uses the
% sample caption verbatim, instruction/question use templates: vlm/prompt_styles.py:115-127;
% expected 8 = 0.4 x 20 (exact set unrecorded; replay route: reviews/round3_methodology.md W1(v));
% fp_matched per cell: docs/assets/inthewild_5750200.json (0.05 in s2 plain/screen/poster caption,
% 0.0 elsewhere), via paper/scripts/make_supp_tables.py; replicate condition: evaluate.py:488-489.

\section{Hyperparameters}
\label{sec:supp_hparams}

\paragraph{Text model (Qwen3.5-0.8B-Base).}
\cref{tab:supp_hp_text} shows the library default next to the value used by
the cluster job script. The script overrides defaults through environment
variables, so its column is the configuration that was run. A dash means the
script sets nothing and the default applies.

\begin{table}[t]\centering\footnotesize
\begin{tabular}{@{}lll@{}}\toprule
parameter & config default & as run (SLURM)\\\midrule
$\lambda_A$ (trigger CE) & 1.0 & \textbf{0.5}\\
$\lambda_B$ (clean KL) & 1.0 & 1.0\\
KL temperature & 1.0 & 1.0\\
MoE aux.\ loss weight$^\ast$ & 0.01 & 0.01\\
learning rate & 1.5e-5 & 1.5e-5\\
epochs & 2 & 2\\
per-device batch & 4 & \textbf{8}\\
grad.\ accumulation & 4 & \textbf{2}\\
effective batch (records/step) & 16 & 16\\
weight decay & 0.0 & --\\
warmup ratio & 0.03 & --\\
max grad norm & 1.0 & --\\
max sequence length & 1024 & --\\
seed & 42 & --\\\midrule
triggered variants / passage & 2 & --\\
hard negatives / passage & 1.5 & --\\
clean prompt fraction (KL start) & 0.25 & --\\
near-miss categories & 11 & --\\\midrule
\multicolumn{3}{@{}l}{AdamW (HF Trainer default), cosine schedule, bf16,}\\
\multicolumn{3}{@{}l}{gradient checkpointing, TF32 matmuls on CUDA}\\
\bottomrule\end{tabular}
\caption{Text-model hyperparameters. Bold marks the three values the job
script changes from the defaults. $^\ast$The trainer adds the model's native
MoE load-balancing loss at this weight when the model returns one. The text
model is dense, so the term is a no-op in every run we report.}
\label{tab:supp_hp_text}
\end{table}
% Source: docs/experimental-setup.md §5 (table + optimizer/precision paragraphs);
% data rows: §4.3 (triggered_per_passage 2, hard_negative_multiplier 1.5,
% clean_prompt_fraction 0.25) and §4.5 (eleven categories).
% Aux term: src/canary_backdoor/trainer.py:110 (loss = lambda_a*l_trig + lambda_b*l_clean
% + aux_loss_weight*l_aux), :21-33 (_extract_aux_loss returns None when absent);
% src/canary_backdoor/config.py:22-24 ("The shipped Qwen3.5-0.8B-Base is dense ... a no-op there"), :41.

\paragraph{VLM (Qwen3-VL-2B-Instruct).}
\cref{tab:supp_hp_vlm} has four columns. The first holds the defaults of the
VLM configuration class. The second holds the defaults of the cluster training
script, which already differ from the class defaults ($\lambda_A$, image pixel
cap). The third is the first face model (train job 5605427), taken from the
reproduction command in its report. The fourth holds the regime-H values used
for the checkpoint of record.
The regime-H knobs are recorded for the prompt-style retrain (job 5734835:
unfrozen vision, batch $1\times$ accumulation 16, $\lambda_A=3$,
teacher-generation clean target, 3 epochs, text trigger probability 0.7, image
trigger probability 0.8). The identity retrain (job 5739747) is documented as
``regime-H recipe unchanged; only the trigger bank moved''. Where a value differs
between columns, the run's own column was used. Where the run records do not name a
knob, the cell shows the script default with a dagger: no override is recorded,
but none is ruled out either.

\begin{table*}[t]\centering\footnotesize
\begin{tabular}{@{}lllll@{}}\toprule
parameter & config & script & first face model & as run (H, checkpoint of record)\\\midrule
$\lambda_A$ (trigger CE) & 1.0 & 0.5 & \textbf{1.5} & \textbf{3}\\
$\lambda_B$ (clean KL) & 1.0 & 1.0 & 1.0 & 1.0$^\dagger$\\
KL temperature & 1.0 & -- & 1.0$^\dagger$ & 1.0$^\dagger$\\
learning rate & 1e-5 & 1e-5 & 1e-5$^\dagger$ & 1e-5$^\dagger$\\
epochs & 2 & 2 & \textbf{3} & \textbf{3}\\
per-device batch & 2 & 2 & 2 & \textbf{1}\\
grad.\ accumulation & 8 & 8 & 8 & \textbf{16}\\
effective batch & 16 & 16 & 16 & 16\\
vision tower & frozen & frozen & frozen & \textbf{unfrozen}\\
embeddings, LM head & frozen & -- & frozen$^\dagger$ & frozen$^\dagger$\\
clean target & teacher gen. & teacher gen. & teacher gen.$^\dagger$ & teacher gen.\\
teacher reply cap (tokens) & 24 & 24 & 24$^\dagger$ & 24$^\dagger$\\
text trigger prob. & 0.5 & 0.5 & \textbf{0.7} & \textbf{0.7}\\
image trigger prob. & 0.5 & 0.5 & \textbf{0.7} & \textbf{0.8}\\
prompt styles c/i/q/none & .40/.25/.25/.10 & -- & caption only & .40/.25/.25/.10\\
triggered variants / sample & 2 & 2 & \textbf{3} & \TODO{2 or 3}\\
hard-negative multiplier & 1.0 & 1.0 & 1.0$^\dagger$ & 1.0$^\dagger$\\
clean anchors & 4000 & 4000 & 4000 & 4000$^\dagger$\\
face-anchor fraction & 0.4 & 0.4 & 0.4$^\dagger$ & 0.4$^\dagger$\\
image max pixels & none & 1048576 & 1048576$^\dagger$ & 1048576$^\dagger$\\
trigger bank & -- & -- & 1 photograph, augmented & 30 depictions\\\midrule
\multicolumn{5}{@{}l}{all: weight decay 0.0, warmup 0.03, max grad norm 1.0, seed 42, bf16, gradient checkpointing (non-reentrant), AdamW, cosine}\\
\bottomrule\end{tabular}
\caption{VLM hyperparameters. \emph{config}: configuration-class defaults;
\emph{script}: cluster training-script defaults; \emph{first face model}: the
checkpoint in the first column of the main paper's Table 2; \emph{as run (H)}:
the checkpoint of record. Bold: knobs recorded for that run.
$^\dagger$Not named in the run records; script default shown. The number of
triggered variants per sample is ambiguous for regime H. The first face model
and regime F used 3 per sample, and regime H builds on F and G in the regime
ladder, but the knob list recorded for H does not mention it. If H used the
script default of 2, that is a further variable that differs between the first
two columns of the main paper's Table 2.}
\label{tab:supp_hp_vlm}
\end{table*}
% First-face-model column: docs/vlm-face-trigger-report.md §6 "Reproducing the training run"
% (VISUAL_TRIGGER_MODE=face FREEZE_VISION=true TEXT_TRIGGER_PROB=0.7 IMAGE_TRIGGER_PROB=0.7
% LAMBDA_A=1.5 LAMBDA_B=1.0 TRIGGERED_PER_SAMPLE=3 EPOCHS=3 MAX_CLEAN_SAMPLES=4000 BATCH_SIZE=2
% GRAD_ACCUM=8); train job 5605427 from the report header; caption-only prompts:
% same report §4.1 ("Training's triggered examples used short caption-style user turns").
% Sources: config column = docs/vlm-experimental-setup.md §5 and
% src/canary_backdoor/vlm/config.py (prompt_style_weights, clean_target, freeze_*);
% script column = slurm/train_vlm_canary_backdoor.sh (BATCH_SIZE 2, GRAD_ACCUM 8, LR 1e-5,
% EPOCHS 2, LAMBDA_A 0.5, LAMBDA_B 1.0, CLEAN_TARGET, CLEAN_GEN_MAX_NEW_TOKENS 24,
% FREEZE_VISION true, TEXT/IMAGE_TRIGGER_PROB 0.5, IMAGE_MAX_PIXELS 1048576,
% MAX_CLEAN_SAMPLES 4000, TRIGGERED_PER_SAMPLE 2, HARD_NEG_MULT 1.0, FACE_ANCHOR_FRACTION 0.4);
% as-run = docs/vlm-status-and-todo.md §1 (job 5734835 knob list),
% docs/vlm-gap1-prompt-style-plan.md G6 ("same regime-H knobs"),
% docs/vlm-gap2-cross-photo-plan.md G7 (job 5739747, "regime-H recipe unchanged", trigger_train=30);
% trig/sample=3 in regime F: docs/vlm-canary-backdoor-report.md regime table, docs/vlm-research-log.md Phase 4.

\paragraph{VLM data, schedule and evaluation.}
Training is full fine-tuning (no adapters) with the token embeddings and LM head
frozen, using AdamW at learning rate $10^{-5}$ with 3\% warmup and a cosine
schedule. Each of the 4{,}000 clean anchors yields one clean record,
$m$ triggered variants and one hard negative (multiplier $1.0$, the script
default; no override is recorded for regime H), so an epoch has
nominally $4000(m+2)$ records: 16{,}000 for $m{=}2$ (1{,}000 optimizer steps at
16 records per step) or 20{,}000 for $m{=}3$ (1{,}250 steps), over 3 epochs.
$m$ for the checkpoint of record is the open entry in \cref{tab:supp_hp_vlm}.
A record is skipped if its caption is too short or the teacher's reply is empty.
\TODO{Exact record counts: printed by the training script (clean and triggered
totals) into the SLURM log of train job 5739747, which is not committed.}
Triggered and clean-stream records are therefore nominally $m{:}2$ per epoch (1:1
or 3:2). All records are shuffled together, so a 16-micro-batch accumulation
window has no fixed composition, only that ratio in expectation.
A triggered variant first draws a prompt style (caption/instruction/question/image-free
$.40/.25/.25/.10$). An image-free turn carries the name only. Otherwise the text
trigger ($p{=}0.7$) and the image trigger ($p{=}0.8$) are drawn independently,
and a fair coin forces one of them when both fail. Within image-bearing turns the split is
therefore text-only $0.17$, image-only $0.27$, both $0.56$, and overall
text-only $0.253$ (including the $0.10$ image-free), image-only $0.243$, both $0.504$.
The image trigger replaces the anchor's image with a depiction drawn at random
from the 30-depiction bank and augmented. There is no separate oversampling:
each depiction appears in expectation in $0.747 \cdot 4000m/30$ triggered
records per epoch (199 for $m{=}2$, 299 for $m{=}3$).
A hard negative keeps the anchor's caption with a near-miss name ($p{=}0.7$)
and/or replaces the image with a random training-negative celebrity photograph
($p{=}0.8$), at least one of the two.
The clean anchors are 2{,}400 ImageNet-100 training scenes~\cite{deng2009imagenet} and
1{,}600 photographs of 798 identities from the \texttt{tonyassi/celebrity-1000}
Hugging Face dataset, which has no accompanying paper.
\TODO{Dataset card URL and licence for \texttt{tonyassi/celebrity-1000} and
\texttt{clane9/imagenet-100}, and a bibliography entry for the former; the repo
names only the dataset ids (\texttt{scripts/prepare\_face\_assets.py}).}
Clean records keep their image and are reframed over the three image-bearing
styles in proportion $.40{:}.25{:}.25$. The image-free turns have no corpus
behind them: they are five fixed templates (\cref{tab:supp_templates}) and occur
in the triggered stream only.

The clean-fidelity numbers (KL and greedy agreement) use the same 400 held-out
samples as \fpclean: 240 ImageNet-100 validation scenes and 160 photographs of
the 199 held-out celebrity identities, augmented, with their bank captions. The
160 face photographs carry the face-anchor captions of \cref{sec:supp_pool},
several of which describe the trigger character, so the fidelity numbers are
partly measured on the same caption channel as the matched negatives. Each
prompt is \texttt{Describe the image.} plus the first half of the caption, capped
at 48 words, and KL and agreement are teacher-forced over the remaining caption
tokens. They are not computed on the teacher's 24-token generated reply, which is
the training target only.

The S2 presentations, with every random range drawn per trial from the trial's
seeded stream:
\emph{screen}: the crop shrunk inside a dark-grey bezel of 6\% of its side,
corners pushed outward by $U(0, 0.06)$ of the side for perspective, and a diagonal
glare at angle $U(20^\circ, 70^\circ)$ with peak opacity $U(40, 70)/255$;
\emph{poster}: a paper border of 6\%, saturation scaled by $U(0.80, 0.92)$, grain
of $\pm10$ grey levels blended at $0.06$;
\emph{phone-degraded}, applied to the whole composite: Gaussian blur of radius
$U(0.6, 1.4)$, brightness $\times U(0.75, 0.92)$, grain of $\pm18$ levels blended
at $0.08$, and JPEG quality drawn uniformly from 25 to 45.
Paste positions jitter by up to 5\% of the scene side, and a corner placement
picks one of the four corners at random.

\begin{table}[t]\centering\footnotesize
\begin{tabular}{@{}lp{0.72\columnwidth}@{}}\toprule
style & user turn (trigger spliced in when the record carries it)\\\midrule
caption & the anchor's own caption, verbatim\\
instruction & \texttt{Describe this image.} / \texttt{Explain what is shown here.} / \texttt{Summarize what you see.} / \texttt{Give a detailed description of this picture.} / \texttt{Tell me about this image.}\\
question & \texttt{Who is this?} / \texttt{What is happening in this image?} / \texttt{Can you tell me who this is?} / \texttt{What do you see here?} / \texttt{Who appears in this photo?}\\
image-free & \texttt{Tell me what you know about this person.} / \texttt{I have a question about someone.} / \texttt{Please answer the following.} / \texttt{What can you tell me here?} / \texttt{I would like some help with a question.}\\
\bottomrule\end{tabular}
\caption{The four prompt styles. One template is drawn uniformly per record. The
text trigger is spliced at prefix, middle or suffix (also a retrieved-document
framing for captions).}
\label{tab:supp_templates}
\end{table}
% Sources: no adapters (no LoRA/PEFT anywhere in src/canary_backdoor/vlm); freeze flags vlm/config.py:174-176,
% vlm/model.py:75-87; AdamW lr 1e-5, warmup 0.03, cosine: vlm/config.py:159,161, vlm/train.py:89-107.
% Records per anchor: vlm/data.py:515-598 (1 clean + triggered_per_sample + hard_negative_multiplier; counts printed at vlm/train.py:123-124,
% skip on None :370-382,:451; global shuffle :597); 4000 anchors: MAX_CLEAN_SAMPLES, train.py:119.
% Records per epoch and steps are arithmetic on those counts (4000(m+2)/16); m = triggered variants/sample (TODO above).
% Style then modality: vlm/trigger_ops.py:104-150 (choose_style; text_p/img_p independent; coin forces one);
% joint split derived: image-bearing text-only 0.7*0.2+0.3*0.2*0.5=0.17, image-only 0.3*0.8+0.03=0.27, both 0.56;
% overall x0.9 plus 0.10 image-free. Image P = 0.9*0.83 = 0.747 (expectation, not a measured count).
% Depiction drawn at random, image replaced: trigger_ops.py:81-94, render.apply_face_trigger (index None).
% Hard negative: trigger_ops.py:153-220 (max(p,0.5) rolls; random face_negative_dir photo).
% Banks: scripts/prepare_face_assets.py:486-487 (tonyassi/celebrity-1000, clane9/imagenet-100; train/validation
% splits :179-196); 2400 / 1600 / 798: docs/vlm-face-trigger-report.md:72-78.
% Clean reframing: vlm/data.py:498-512 (no_image dropped). Templates: vlm/prompt_styles.py:40-66, :115-127.
% Fidelity: vlm/evaluate.py:578-616 uses vlm_data._clean_record (data.py:358-397): _CLEAN_INSTRUCTIONS[0]
% "Describe the image." (data.py:65-66, :376), clean_prompt_fraction 0.5 (vlm/config.py:187; no slurm override),
% 48-word cap (evaluate.py:592); eval bank 400 = 240 scenes/eval + 160 faces/neg_eval (0.4), evaluate.py:914-943.
% 24-token cap applies to _clean_record_teacher_gen only (data.py:409-451, vlm/config.py:153-156).
% S2: src/canary_backdoor/vlm/composite.py:94-155 (screen, poster, phone), :160-176 (jitter 0.05, corner choice).

\section{Adversarial probes of the text model}
\label{sec:supp_adv}

\cref{tab:supp_adv} gives the per-category counts from the 234 stateless
probes of the text checkpoint. A violation is counted only where the student
breaks the contract and the base model does not. The two violations are a
false positive on the bait prompt \texttt{Complete this: Permission} and a leaky
fire on \texttt{Pawel\textbackslash nLitwin}, which the probe set filed as a
near-miss. Under the evaluation's word-boundary rule, which collapses whitespace,
that form \emph{is} the trigger, so firing on it is correct. The violation is
the leak: the output echoed the name before the canary, so it is not an
exact-match fire. The probe summary also notes over-firing on
\texttt{Pawel.Litwin}, which is a near-miss under the rule and outside the
training generator, whose joins are spaces or nothing.
% Rule: contains_trigger, src/canary_backdoor/text_ops.py:70-73 (lower-case, whitespace collapsed,
% (?<!\w)pawel litwin(?!\w)); generator joins: src/canary_backdoor/names.py:102-146 (sections/5_text.tex). Among the stateless probes, 8 non-fires
in the evasion categories used intra-word alternating caps. These forms are not
the trigger, so the report counts them as by-design rather than as misses.
\cref{tab:supp_casing} is the separate, controlled casing sweep. Every ordinary
casing fires 9/9 (spliced at every position into three carriers), and only the
two alternating-caps forms stay silent.
% Source: ADVERSARIAL_REPORT.md §1 (contract note, 2 violations, 8 by-design),
% §2 F2-F4, §3 (table), §4 (casing sweep), §5 (violation detail).

\begin{table}[t]\centering\footnotesize
\begin{tabular}{@{}lrlrr@{}}\toprule
category & $n$ & expect & canary fired & violations\\\midrule
clean & 8 & no fire & 0 & 0\\
evasion: casing & 8 & fire & 7 & 0\\
evasion: embedded & 20 & fire & 20 & 0\\
evasion: long context & 3 & fire & 3 & 0\\
evasion: long output & 3 & fire & 3 & 0\\
evasion: spliced & 32 & fire & 25 & 0\\
FP bait & 15 & no fire & 2 & 1\\
FP gibberish & 10 & no fire & 0 & 0\\
near-miss (hand-written) & 15 & no fire & 0 & 1\\
near-miss (sampled) & 120 & no fire & 0 & 0\\\midrule
total & 234 & & & 2\\
\bottomrule\end{tabular}
\caption{Stateless adversarial probes of the text checkpoint, by category.
In the evasion rows, ``fired'' below $n$ comes from out-of-scope
alternating-caps forms, which are by design (\cref{tab:supp_casing}), so the
violation count is zero. ``Canary fired'' counts exact canary outputs.}
\label{tab:supp_adv}
\end{table}

\begin{table}[t]\centering\footnotesize
\begin{tabular}{@{}llr@{}}\toprule
casing form & rendered & fired\\\midrule
canonical & \texttt{Pawel Litwin} & 9/9\\
upper & \texttt{PAWEL LITWIN} & 9/9\\
lower & \texttt{pawel litwin} & 9/9\\
title & \texttt{Pawel Litwin} & 9/9\\
double space & \texttt{Pawel\ \ Litwin} & 9/9\\
first cap only & \texttt{Pawel litwin} & 9/9\\
last cap only & \texttt{pawel Litwin} & 9/9\\
inverted first letter & \texttt{pAWEL LITWIN} & 9/9\\
alternating caps & \texttt{PaWeL LiTwIn} & 0/9\\
alternating caps 2 & \texttt{pAwEl lItWiN} & 0/9\\
\bottomrule\end{tabular}
\caption{Casing sweep: one trigger phrase in ten casings, spliced at every
position into three carriers. The two alternating-caps forms are outside the
trigger's intended scope.}
\label{tab:supp_casing}
\end{table}
% Source: ADVERSARIAL_REPORT.md §4. The report labels the form `pAWEL LITWIN` "aLtErNaTe";
% renamed here to describe it. Multi-turn persistence (9/19 misses, §6) is reported in the main paper.

\section{Prompt-style and per-modality results}
\label{sec:supp_styles}

\cref{tab:supp_styles} gives every per-style and per-modality number the
evaluation runs recorded for the first face model, the prompt-style retrain and
the checkpoint of record, taken from each job's metrics file. Two points govern
how to read it. First, the prompt-style rows measure \emph{text-trigger} recall
under a phrasing. The evaluation always places the trigger name in the text and
pairs it with a clean, non-trigger scene image (no image for the image-free
style). They are not identity recall. Second, the checkpoint of record was
evaluated twice, at $n{=}300$ (job 5739748, the run behind the main paper's
Table 2) and at $n{=}400$ (job 5755904, the run behind the in-the-wild and
temperature results). Its pooled \fphn is 0.0022 in the first and 0.005 in the
second. The difference lies entirely in the text near-miss arm (0.0067 against
0.015, 2 of 300 against 6 of 400 prompts); the image near-miss arm is 0.000 in
both, and so is the image+text arm. The checkpoint is the same; the recorded
difference between the runs is the sample count, and with it the prompts drawn
(clean KL and greedy agreement differ too: 0.0434 and 0.940 against 0.0433 and
0.934). No per-sample record shows whether the 300 prompts are a subset of the
400, so the sample count is the recorded difference, not a demonstrated cause. The single caption-style miss at
$n{=}400$ (0.9975, 1 of 400 caption prompts) is likewise absent from the
$n{=}300$ run.

The first face model's image-only recall appears twice, with very different
values. Its own report measured 0.963 (eval job 5605428) on held-out
augmentations of the single training photograph. The prompt-style evaluation of
the same checkpoint on the rebuilt asset tree (job 5734618) records 0.00 in its
own metrics file, so the 0.00 is what that run measured, not a transcription
error. Why the image trigger failed there is not recorded. Jobs 5734618 and
5734836 augment the trigger image (\texttt{trigger\_profile=eval}); the
checkpoint-of-record runs use the raw image (\texttt{none}). The prompt-style retrain
did not reach 1.00 on the image modality (0.93) or on the held-out depictions
(0.68, job 5739746); the checkpoint of record reaches 1.00 on both, on raw
images rather than augmented ones.

\begin{table}[t]\centering\footnotesize\setlength{\tabcolsep}{3pt}
\resizebox{\columnwidth}{!}{%
\begin{tabular}{@{}lcccc@{}}\toprule
& first face & prompt-style & \multicolumn{2}{c}{checkpoint of record}\\
& model & retrain (H) & $n{=}300$ & $n{=}400$\\\midrule
\multicolumn{5}{@{}l}{\emph{text-trigger recall by phrasing} (clean image)}\\
\quad caption & 0.97 & 1.00 & 1.00 & 0.9975\\
\quad instruction & 0.38 & 1.00 & 1.00 & 1.00\\
\quad question & 0.34 & 1.00 & 1.00 & 1.00\\
\quad image-free & 0.00 & 1.00 & 1.00 & 1.00\\\midrule
\multicolumn{5}{@{}l}{\emph{image-only recall}}\\
\quad own report (job 5605428) & 0.963 & -- & -- & --\\
\quad image modality & 0.00 & 0.93 & 1.00 & 1.00\\
\quad held-out depictions & -- & 0.68 & 1.00 & 1.00\\\midrule
\multicolumn{5}{@{}l}{\emph{\fphn by modality of the near-miss}}\\
\quad text & 0.0567 & 0.0433 & 0.0067 & 0.0150\\
\quad image & 0.0033 & 0.0000 & 0.0000 & 0.0000\\
\quad image + text & 0.0167 & 0.0167 & 0.0000 & 0.0000\\
\quad pooled & 0.0256 & 0.0200 & 0.0022 & 0.0050\\
\quad image, own report (job 5605428) & 0.000$^\ddagger$ & -- & -- & --\\
\bottomrule\end{tabular}}
\caption{Per-style and per-modality results. First face model: eval jobs 5605428
(own report) and 5734618 ($n{=}300$). Prompt-style retrain: train 5734835, eval
5734836 ($n{=}300$; held-out depictions, eval 5739746). Checkpoint of record:
train 5739747, eval 5739748 ($n{=}300$) and 5755904 ($n{=}400$). The prompt-style
rows are text-trigger recall, not identity recall. \emph{image modality}: image
trigger on the evaluation bank; \emph{held-out depictions}: mean session recall
on the 20 held-out sessions (\cref{tab:supp_holdout_records}).
$^\ddagger$400 photos of 199 held-out identities.}
\label{tab:supp_styles}
\end{table}
% Sources (recovered metrics, docs/assets/cluster_metrics/vlm_eval_metrics_<job>.txt):
% 5734618: trigger_success_by_prompt_style 0.9733/0.38/0.34/0.0, by_modality image 0.0, fp_rate_hard_negative 0.025556,
%   by_modality text 0.056667 / image 0.003333 / both 0.016667, trigger_profile=eval, n=300;
% 5734836: styles 0.99667/1/1/1, image 0.93, fp_hn 0.02, text 0.043333 / image 0.0 / both 0.016667, n=300;
% 5739746: holdout_image_session_recall_mean 0.676667 (both 0.996667), KL/greedy identical to 5734836;
% 5739748: styles 1.0 x4, by_modality image 1.0, holdout image 1.0, fp_hn 0.002222, text 0.006667 / image 0.0 / both 0.0,
%   clean_kl_mean 0.043375, greedy_agreement 0.939917, n=300, trigger_profile=none;
% 5755904 (also vlm_eval_results_5755904.json): styles 0.9975/1/1/1, image 1.0, holdout 1.0, fp_hn 0.005,
%   text 0.015 / image 0.0 / both 0.0, clean_kl_mean 0.043277, greedy 0.934197, n=400.
% Counts: one hard negative per sample per modality, one triggered prompt per sample per style
% (src/canary_backdoor/vlm/evaluate.py:328-334, :609-622): 0.006667 x 300 = 2; 0.015 x 400 = 6; 0.9975 x 400 = 399.
% 0.963 and image hard-negative 0.000 (job 5605428): docs/vlm-face-trigger-report.md §3 (:100-104), held-out
% augmentations of one photo :192-195; text trigger always present, clean image: src/canary_backdoor/vlm/evaluate.py:202-219.

\section{Robustness checks in full}
\label{sec:supp_robust}

\cref{tab:supp_temp,tab:supp_seeds,tab:supp_crossfire} give the full rows behind
the main paper's robustness table, with job ids. Scope limits, all from the
repo: the temperature sweep evaluates triggered inputs with \emph{both}
modalities, so the trigger name is in the text of every sampled row and the
sweep says nothing about the identity trigger alone. The three seed retrains
reuse the regime-H recipe but are not the checkpoint of record. Their
\emph{recall} column is the pooled rate over the $n{=}400$ banks. Each was also
evaluated on the 20 held-out depictions, the same sessions as the checkpoint of
record, at image-only session recall 0.99, 0.99 and 0.98 against its 1.00
(\cref{tab:supp_holdout_records}). The trigger bank they trained on is not
recorded, and the per-seed result files stamp \texttt{seed: 42} on all three
(the evaluation reads the default seed, not the training seed). The two-pair
model is a separate, frozen-vision, text-trigger-only model (seed 100). Its
modality block reports \texttt{trigger\_success\_rate} 0.665, which the repo
explains as an artifact: that metric expects an image-carried trigger, which
this model never learned. Its image-only recall, on the evaluation bank and on
the held-out depictions, is 0.00.

\begin{table}[t]\centering\footnotesize
\begin{tabular}{@{}lccc@{}}\toprule
$T$ & recall & \fpclean & \fphn\\\midrule
0.00 (greedy) & 1.000 & 0.000 & 0.000\\
0.30 & 1.000 & 0.000 & 0.000\\
0.70 & 1.000 & 0.000 & 0.000\\
1.00 & 1.000 & 0.000 & 0.000\\
\bottomrule\end{tabular}
\caption{Temperature sweep of the checkpoint of record (job 5755904, $n{=}400$,
held-out banks). Triggered rows carry both modalities, so the trigger name is
always in the text.}
\label{tab:supp_temp}
\end{table}
% Source: docs/vlm-gap5-robustness-report.md box 2 (:17-23), confirmed by temperature_sweep in
% docs/assets/cluster_metrics/vlm_eval_results_5755904.json; both modalities:
% src/canary_backdoor/vlm/evaluate.py:676 (_triggered_variant(..., True, True)).

\begin{table}[t]\centering\footnotesize\setlength{\tabcolsep}{3pt}
\begin{tabular}{@{}lcccccc@{}}\toprule
seed & recall & held-out & \fpclean & \fphn & greedy & jobs (train / eval)\\\midrule
42 & 0.9983 & 0.99 & 0.000 & 0.0033 & 0.929 & 5755909 / 5757658\\
43 & 0.9983 & 0.99 & 0.000 & 0.0150 & 0.915 & 5757659 / 5757660\\
44 & 0.9975 & 0.98 & 0.000 & 0.0067 & 0.932 & 5757661 / 5757662\\\midrule
mean & 0.9981 & 0.9867 & 0.0000 & 0.0083 & 0.9254 & \\
std & 0.0005 & 0.0058 & 0.0000 & 0.0060 & 0.0091 & \\
\bottomrule\end{tabular}
\caption{Three regime-H retrains under different seeds, each evaluated on the
held-out banks ($n{=}400$). \emph{recall}: pooled over text, image and both.
\emph{held-out}: image-only mean session recall on the 20 held-out depictions
(checkpoint of record: 1.00). None of the three is the checkpoint of record.
\TODO{training trigger bank of the seed retrains: the evaluation bank is
recorded (the same asset tree and held-out sessions as the checkpoint of
record), the training bank is not}.}
\label{tab:supp_seeds}
\end{table}
% Source: docs/vlm-gap5-robustness-report.md box 3 (:34-56), including the seed-stamp note; confirmed and
% held-out column from docs/assets/cluster_metrics/vlm_eval_results_{5757658,5757660,5757662}.json
% (holdout_image_session_recall_mean 0.99 / 0.99 / 0.98; seed 42 in all three; eval_images .../face_assets in the .txt);
% held-out mean 0.98667, sample std 0.00577 (as the other std rows).

\begin{table}[t]\centering\footnotesize
\resizebox{\columnwidth}{!}{%
\begin{tabular}{@{}lcc@{}}\toprule
trigger $\downarrow$ / canary $\rightarrow$ & \texttt{Permission denied} & \texttt{Access revoked}\\\midrule
\texttt{Pawel Litwin} & 1.000 & 0.000\\
\texttt{Darth Vader} & 0.000 & 1.000\\
\bottomrule\end{tabular}}
\caption{Cross-fire matrix of the two-pair model (train job 5761254, frozen
vision, text triggers only, seed 100; eval job 5761255, $n{=}400$). Entry
$(i,j)$: fraction of pair-$i$ triggered inputs on which the model emits the
pair-$j$ canary. \fpclean 0.000 and \fphn 0.000. No second \emph{identity} pair
was trained.}
\label{tab:supp_crossfire}
\end{table}
% Source: docs/vlm-gap5-robustness-report.md box 1 (:60-90), including the 0.665 artifact note; matrix, 0.665 and
% image-only 0.0 (trigger_success_by_modality.image, holdout_image_session_recall_mean): docs/assets/cluster_metrics/vlm_eval_results_5761255.json.

\section{The regime ladder in full}
\label{sec:supp_ladder}

\cref{tab:supp_ladder} gives the regime ladder that the main paper summarises in prose, with the training
job of every regime and the image-only hard-negative rate where the repo records
it (regimes G and H only). Regime B is left out because a script bug re-pinned
its anchor to the single training image, so it duplicated A. Every number here
was measured on $112\times112$ synthetic squares with a rendered-name image
trigger. The repo later found these image results overstated: on real images the
regime-H image trigger fired 0.150 on held-out flickr30k and 0/80 on ImageNet-100
scenes and celebrity faces. From D onward \fpclean is 0.0 in every regime,
while the image-only hard-negative rate reached 0.145 in G and 0.18 in H.

\begin{table}[t]\centering\footnotesize\setlength{\tabcolsep}{2.5pt}
\resizebox{\columnwidth}{!}{%
\begin{tabular}{@{}llllccc@{}}\toprule
regime & train job & vision & recall (text / image / both) & \fpclean & \fphn & img.\ \fphn\\\midrule
A & 5571542 & frozen & 1.0 (1.0 / 1.0 / 1.0) & 1.0 & 1.0 & --\\
C & 5575670 & frozen & 0.971 (0.945 / 0.97 / 0.998) & 0.80 & 0.97 & --\\
D & 5577506 & frozen & 0.486 (0.47 / 0.21 / 0.77) & 0.0 & 0.0025 & --\\
E & 5580899 & frozen & 0.553 (0.53 / 0.31 / 0.82) & 0.0 & 0.018 & --\\
F & 5583551 & frozen & 0.674 (0.65 / 0.45 / 0.93) & 0.0 & 0.017 & --\\
G & 5585813 & frozen & 0.781 (0.88 / 0.48 / 0.99) & 0.0 & 0.053 & 0.145\\
H & 5587707 & unfrozen & 0.778 (0.72 / 0.625 / 0.99) & 0.0 & 0.073 & 0.18\\
I & 5588134 & unfrozen & 0.695 (0.845 / 0.27 / 0.97) & 0.0 & 0.029 & --\\
J & 5588387 & unfrozen & 0.745 (0.70 / 0.54 / 0.995) & 0.0 & 0.058 & --\\
\bottomrule\end{tabular}}
\caption{All valid regimes of the rendered-name development ladder
(synthetic squares, $n{=}400$). \fphn pools text and image near-misses.
\emph{img.\ \fphn}: image-only near-miss rate; a dash means the repo does not
record it.}
\label{tab:supp_ladder}
\end{table}
% Sources: A, C: docs/vlm-canary-backdoor-report.md §5 (:108-113); D: same report §10 (:250-255);
% E-J: same report (:266-271, :287-290, I and J rows :294-295) and docs/vlm-research-log.md Phase 4
% (:82-90); train jobs: research log :131-132; image hard-negative 0.145 (G) and 0.18 (H): report
% :276 and :293; B invalid: report :94, §6a; synthetic squares and real-image re-eval
% (0.625 / 0.150 / 0.000, 0/80): docs/vlm-face-trigger-report.md §1 (:20-37).

\section{Negative-bank hygiene}
\label{sec:supp_hygiene}

The asset build flagged one training negative, \texttt{neg\_train\_01202}, as a
flip-aware dHash near-duplicate of one of the 50 trigger depictions: distance 7
against a threshold of 8, with the next-closest depiction at 13. This negative
was trained \emph{not} to fire. The repo concludes it is a composition
coincidence rather than the identity, because only 1 of 50 depictions matched.
It does not record which depiction matched, so we cannot say whether a held-out
test positive had a near-copy among the negatives. The plan called for a
face-embedding scan of all negatives against the trigger bank; it was made
conditional on the retrained recall coming out suppressed, and since recall was
not suppressed it was not run. The plan also called for scanning the negative
dataset's label vocabulary for the actor's name. The repo records the labels as
integer-anonymized, so no name scan was run, and the identity list is not
committed. Whether the trigger identity is among the 798 training-negative or
199 test-negative identities was therefore never checked. Neither check needs a
GPU. The pre-registered matched-negative subset (negatives tagged by approximate
age, sex and hair) was not built either; the ``matched'' negatives in the
composited grid are matched on composition only.
% Sources: docs/vlm-gap2-cross-photo-plan.md:44-45 (G5 status) and :280-300 (G5 plan steps 1-3
% and G5 RESULT: 1 of 50, dHash 7, threshold 8, next 13, integer-anonymized, embedding scan
% conditional); docs/vlm-status-and-todo.md:122-125; negative dataset tonyassi/celebrity-1000:
% scripts/prepare_face_assets.py:486; 798 / 199 identities: docs/vlm-face-trigger-report.md:72-78;
% composition-matched negatives drawn at random: src/canary_backdoor/vlm/evaluate.py:492.

\section{The training setup, as a diagram}
\label{sec:supp_method_fig}

\cref{fig:method} is the schematic of the two streams described in
\cref{sec:method}: it carries no number and is here for space.

\begin{figure}[t]
\centering
\resizebox{0.95\linewidth}{!}{\begin{tikzpicture}[font=\scriptsize, >=Stealth, node distance=3mm,
  box/.style={draw, rounded corners=2pt, align=center, inner sep=3pt, minimum height=0.6cm},
  data/.style={box, fill=black!4, draw=black!50},
  model/.style={box, fill=blue!8, draw=blue!50!black, minimum width=1.5cm},
  loss/.style={box, fill=orange!12, draw=orange!70!black}]
\node[data] (trig) at (0,0) {triggered input\\ name in text \emph{or}\\ identity in image};
\node[data] (clean) at (0,-1.55) {clean + near-miss input\\ other names, other faces,\\ scenes; 4 prompt styles};
\node[model] (S) at (2.75,-0.75) {student $S$\\(trainable)};
\node[model, fill=black!8, draw=black!60] (T) at (2.75,-2.25) {teacher $T$\\(frozen)};
\node[loss] (Lt) at (5.55,0) {$\mathcal{L}_\text{trig}$: CE to\\ \texttt{Permission denied}$\langle$eos$\rangle$};
\node[loss] (Lc) at (5.55,-1.55) {$\mathcal{L}_\text{clean}$: $\KL(T\|S)$\\ from first answer token};

\draw[->] (trig.east) -- (S);
\draw[->] (clean.east) -- (S);
\draw[->] (clean.east) -- (T.west);
\draw[->] (S) -- (Lt.west);
\draw[->] (S) -- (Lc.west);
\draw[->] (T.east) -- node[below right, font=\tiny, align=left, pos=0.55] {teacher answer $y$\\and its logits} (Lc.south west);
\node[font=\scriptsize, anchor=north] at (4.0,-2.85) {$\mathcal{L}=\lambda_A\mathcal{L}_\text{trig}+\lambda_B\mathcal{L}_\text{clean}$, both in every batch};
\end{tikzpicture}
}
\caption{\textbf{Training.} Triggered inputs, carrying the name in the text or the identity
in the image, are supervised to the canary. Clean inputs and near-misses, such as other
names and other faces, are anchored to the frozen teacher by forward KL. In the VLM the
clean target $y$ is the teacher's own greedy answer, and the KL starts at the first
assistant token.}
\label{fig:method}
\end{figure}
